\section{Search in a Rotated Sorted Array II}
\label{sec:bs-rotated-2}

\problemstatement{Same setting as Section~\ref{sec:bs-rotated-1} -- a
sorted array rotated at an unknown pivot -- except now the array
\textbf{may contain duplicate values}. Given the rotated array
$\textit{arr}$ and a target $X$, determine whether $X$ is present
(\texttt{true}/\texttt{false}).}

\begin{patternbox}
The keyword change from Section~\ref{sec:bs-rotated-1} is just one word:
\emph{``may contain duplicates.''} This single word breaks the comparison
$\textit{arr}[\textit{lo}] \le \textit{arr}[\textit{mid}]$ that we relied on
to decide which half is sorted -- with duplicates, this comparison can be
\texttt{true} even when $\textit{lo}$, $\textit{mid}$, and the rotation
pivot are tangled together in a way that makes \emph{neither} half safely
identifiable as sorted. Whenever a rotated-array problem mentions
duplicates, expect to need an escape hatch for this ambiguous case, and
expect the worst-case time complexity to degrade as a result.
\end{patternbox}

\subsection*{Understanding the problem carefully}

Consider $\textit{arr} = [3,3,3,1,2,3]$ with $\textit{lo}=0$ pointing to
value $3$, and suppose $\textit{mid}$ also happens to land on a $3$ (say
$\textit{mid}=2$, $\textit{arr}[2]=3$). Then
$\textit{arr}[\textit{lo}] = \textit{arr}[\textit{mid}] = 3$, so the
inequality $\textit{arr}[\textit{lo}] \le \textit{arr}[\textit{mid}]$ holds,
which the Section~\ref{sec:bs-rotated-1} logic would interpret as ``the
left half is sorted.'' But the rotation pivot (where $1$ follows $3$)
actually sits \emph{inside} this very range, so the left half is not
genuinely a clean sorted run with respect to deciding where $X$ could be --
the equal values at the boundary hide the pivot's location. With distinct
elements this ambiguity cannot arise (any two unequal boundary values
immediately tell us which side the pivot is on), but duplicates can make
$\textit{arr}[\textit{lo}] = \textit{arr}[\textit{mid}]$ true regardless of
where the pivot actually is.

\begin{ideabox}[{Greedy Decision, with an Escape Hatch}]
Use the exact same logic as Section~\ref{sec:bs-rotated-1}, \emph{except}:
whenever $\textit{arr}[\textit{lo}] = \textit{arr}[\textit{mid}]$ \emph{and}
$\textit{arr}[\textit{mid}] = \textit{arr}[\textit{hi}]$ (or, more simply,
whenever the usual ``which half is sorted'' check is ambiguous because of
equal boundary values), we cannot safely discard either half based on
comparisons alone. In that situation, shrink the window by one element from
each side ($\textit{lo} \gets \textit{lo}+1$, $\textit{hi} \gets
\textit{hi}-1$) and try again -- this never discards a position that could
uniquely hold $X$ in a way that loses correctness, but it does cost us the
usual halving speed in the worst case.
\end{ideabox}

\subsection*{Why shrinking by one is the right (and only) safe fallback}

When $\textit{arr}[\textit{lo}] = \textit{arr}[\textit{mid}]$, we genuinely
cannot tell, from this comparison alone, which half contains the rotation
pivot -- both possibilities are consistent with the observed equality. Since
we cannot identify a provably safe half to discard, we are not entitled to
discard a full half at all without risking skipping over $X$. The one
adjustment we \emph{can} make safely is to shrink the boundary by a single
element: if $\textit{arr}[\textit{lo}]$ duplicates the value at
$\textit{mid}$, removing just $\textit{lo}$ from consideration cannot lose
$X$, because if $X = \textit{arr}[\textit{lo}]$, the duplicate value
persists at $\textit{mid}$ and remains discoverable; if $X \ne
\textit{arr}[\textit{lo}]$, removing $\textit{lo}$ loses nothing relevant.
This is why the worst case degrades to $O(n)$: an adversarial array of
nearly all-equal elements (e.g. $[2,2,2,2,2,2,2,1,2]$) can force this
one-at-a-time shrinking on almost every step.

\subsection*{Algorithm}

\begin{enumerate}
  \item Initialize $\textit{lo} \gets 0$, $\textit{hi} \gets n-1$.
  \item While $\textit{lo} \le \textit{hi}$:
  \begin{enumerate}
    \item $\textit{mid} \gets \textit{lo}+(\textit{hi}-\textit{lo})/2$.
    \item If $\textit{arr}[\textit{mid}] = X$: return \texttt{true}.
    \item If $\textit{arr}[\textit{lo}] = \textit{arr}[\textit{mid}]$ and
          $\textit{arr}[\textit{mid}] = \textit{arr}[\textit{hi}]$:
          $\textit{lo} \gets \textit{lo}+1$, $\textit{hi} \gets
          \textit{hi}-1$ (ambiguous; shrink by one from each side).
    \item Else if $\textit{arr}[\textit{lo}] \le \textit{arr}[\textit{mid}]$
          (left half safely sorted): apply the same range check as
          Section~\ref{sec:bs-rotated-1}.
    \item Else (right half safely sorted): apply the same range check as
          Section~\ref{sec:bs-rotated-1}.
  \end{enumerate}
  \item If the loop ends without returning, return \texttt{false}.
\end{enumerate}

\begin{algorithm}[H]
\caption{Search in Rotated Sorted Array (With Duplicates)}
\begin{algorithmic}[1]
\Procedure{SearchRotatedDup}{$\textit{arr}, X$}
  \State $\textit{lo} \gets 0$, $\textit{hi} \gets n-1$
  \While{$\textit{lo} \le \textit{hi}$}
    \State $\textit{mid} \gets \textit{lo}+(\textit{hi}-\textit{lo})/2$
    \If{$\textit{arr}[\textit{mid}] = X$} \Return \textbf{true} \EndIf
    \If{$\textit{arr}[\textit{lo}] = \textit{arr}[\textit{mid}]$ \textbf{and} $\textit{arr}[\textit{mid}] = \textit{arr}[\textit{hi}]$}
      \State $\textit{lo} \gets \textit{lo}+1$, $\textit{hi} \gets \textit{hi}-1$
    \ElsIf{$\textit{arr}[\textit{lo}] \le \textit{arr}[\textit{mid}]$}
      \If{$\textit{arr}[\textit{lo}] \le X$ \textbf{and} $X < \textit{arr}[\textit{mid}]$}
        \State $\textit{hi} \gets \textit{mid} - 1$
      \Else
        \State $\textit{lo} \gets \textit{mid} + 1$
      \EndIf
    \Else
      \If{$\textit{arr}[\textit{mid}] < X$ \textbf{and} $X \le \textit{arr}[\textit{hi}]$}
        \State $\textit{lo} \gets \textit{mid} + 1$
      \Else
        \State $\textit{hi} \gets \textit{mid} - 1$
      \EndIf
    \EndIf
  \EndWhile
  \State \Return \textbf{false}
\EndProcedure
\end{algorithmic}
\end{algorithm}

\subsection*{Dry run}

\dryrun{Take $\textit{arr} = [3,1,2,3,3,3,3]$, $X = 1$.}

\begin{itemize}
  \item $\textit{lo}=0,\textit{hi}=6$: $\textit{mid}=3$,
        $\textit{arr}[3]=3 \ne 1$. Is $\textit{arr}[0]=3=\textit{arr}[3]=3$
        and $\textit{arr}[3]=3=\textit{arr}[6]=3$? Yes -- ambiguous. Shrink:
        $\textit{lo}=1$, $\textit{hi}=5$.
  \item $\textit{lo}=1,\textit{hi}=5$: $\textit{mid}=3$,
        $\textit{arr}[3]=3 \ne 1$. Is $\textit{arr}[1]=1=\textit{arr}[3]=3$?
        No. Is $\textit{arr}[1]=1 \le \textit{arr}[3]=3$? Yes, left half
        sorted. Is $1 \le 1 < 3$? Yes. Search left: $\textit{hi}=2$.
  \item $\textit{lo}=1,\textit{hi}=2$: $\textit{mid}=1$,
        $\textit{arr}[1]=1=1$. Return \texttt{true}.
\end{itemize}

The answer is \texttt{true}.

\subsection*{C++ implementation}

\begin{lstlisting}
#include <bits/stdc++.h>
using namespace std;

bool searchRotatedDup(vector<int>& arr, int X) {
    int n = arr.size();
    int lo = 0, hi = n - 1;

    while (lo <= hi) {
        int mid = lo + (hi - lo) / 2;
        if (arr[mid] == X) return true;

        if (arr[lo] == arr[mid] && arr[mid] == arr[hi]) {
            lo++;
            hi--;
        } else if (arr[lo] <= arr[mid]) {
            if (arr[lo] <= X && X < arr[mid]) {
                hi = mid - 1;
            } else {
                lo = mid + 1;
            }
        } else {
            if (arr[mid] < X && X <= arr[hi]) {
                lo = mid + 1;
            } else {
                hi = mid - 1;
            }
        }
    }
    return false;
}

int main() {
    vector<int> arr = {3, 1, 2, 3, 3, 3, 3};
    cout << (searchRotatedDup(arr, 1) ? "true" : "false") << endl;
    return 0;
}
\end{lstlisting}

\begin{complexitybox}
\textbf{Time Complexity.} In the best/average case, the same halving as
Section~\ref{sec:bs-rotated-1} applies, giving $O(\log n)$. But in the
worst case -- an array of almost entirely duplicate values, e.g.
$[2,2,2,2,2,2,1,2]$ -- the ambiguous case forces shrinking by one element
at a time, giving a worst-case time complexity of
\[
  O(n).
\]
\textbf{Space Complexity.} $O(1)$ auxiliary space.
\end{complexitybox}

\begin{takeawaybox}
Duplicates are binary search's natural enemy: they tend to destroy the
clean monotonicity or structural invariants that comparisons rely on,
forcing a fallback to linear-time shrinking in the worst case. Whenever a
sheet problem adds ``with duplicates'' to a previously clean binary search
problem, expect to need an explicit ambiguous-case branch, and expect to
state a degraded worst-case complexity honestly rather than claiming
$O(\log n)$ that the duplicates actually prevent.
\end{takeawaybox}
